\begin{frame}{Row-operation example: arithmetic made visible}
  \centering
  \Large
  \[
  \underbrace{(2,4,3)}_{R_2}
  -2\underbrace{(1,2,-1)}_{R_1}
  =
  (2-2,\;4-4,\;3+2)
  =
  \underbrace{(0,0,5)}_{R_2^{\text{new}}}
  \]

  \vspace{0.7cm}
  \begin{columns}[T,onlytextwidth]
    \begin{column}{0.47\textwidth}
      \begin{block}{Cancelled}
        The repeated contribution in columns 1 and 2.
      \end{block}
    \end{column}
    \begin{column}{0.47\textwidth}
      \begin{exampleblock}{Exposed}
        A remaining contribution in column 3.
      \end{exampleblock}
    \end{column}
  \end{columns}
\end{frame}

\begin{frame}{Row-operation example: before, operation, after}
  \begin{columns}[c,onlytextwidth]
    \begin{column}{0.35\textwidth}
      \centering
      \mutedtext{before}
      \[
      \begin{pmatrix}
      1&2&-1\\
      \cellcolor{softorange}2&\cellcolor{softorange}4&\cellcolor{softorange}3\\
      0&1&5
      \end{pmatrix}
      \]
    \end{column}
    \begin{column}{0.25\textwidth}
      \centering
      \pill{softred}{$R_2\leftarrow R_2-2R_1$}\\[0.35cm]
      \Huge $\Longrightarrow$
    \end{column}
    \begin{column}{0.35\textwidth}
      \centering
      \mutedtext{after}
      \[
      \begin{pmatrix}
      1&2&-1\\
      \cellcolor{softgreen}0&\cellcolor{softgreen}0&\cellcolor{softgreen}5\\
      0&1&5
      \end{pmatrix}
      \]
    \end{column}
  \end{columns}

  \vspace{0.45cm}
  \takeaway{The new row shows exactly what was not already explained by the first row.}
\end{frame}

\begin{frame}{The three layers of the same step}
  \begin{columns}[T,onlytextwidth]
    \begin{column}{0.31\textwidth}
      \begin{block}{1. Arithmetic}
        subtract two row vectors
      \end{block}
    \end{column}
    \begin{column}{0.31\textwidth}
      \begin{exampleblock}{2. Transformation}
        replace $R_2$ by a new row
      \end{exampleblock}
    \end{column}
    \begin{column}{0.31\textwidth}
      \begin{alertblock}{3. Interpretation}
        reveal dependence and residual information
      \end{alertblock}
    \end{column}
  \end{columns}

  \vspace{0.8cm}
  \concept{A good solution makes all three layers visible.}
\end{frame}

\begin{frame}{Column vectors are matrices}
  \begin{columns}[T,onlytextwidth]
    \begin{column}{0.44\textwidth}
      \begin{block}{Column form}
        \[
        \bm{u}=\begin{pmatrix}2\\-1\\4\end{pmatrix}
        \in\R^{3\times1}
        \]
      \end{block}
    \end{column}
    \begin{column}{0.44\textwidth}
      \begin{exampleblock}{Row form}
        \[
        \bm{u}^{\mathsf T}=\begin{pmatrix}2&-1&4\end{pmatrix}
        \in\R^{1\times3}
        \]
      \end{exampleblock}
    \end{column}
  \end{columns}

  \vspace{0.55cm}
  \takeaway{Transpose changes the algebraic role of the same coordinates.}
\end{frame}
